% !TeX root = model_v2_standalone.tex
% model_v2.tex
% Category-learning state model v2：精简后的正式主模型方法底稿。
% 本文件只描述主模型中保留的机制，不收录旧版架构搜索中的可选扩展。
% 在正式投稿前，所有尚未冻结的数值搜索边界与模拟次数应与最终代码配置一致。

\subsubsection{模型目标与理论承诺}

本研究使用逐试次计算模型刻画类别学习过程中内部规则信念的形成、维持与修正。模型的核心假设是，实验任务定义了一个完整的候选假设空间 $\mathcal H$，但学习者在任一试次只能同时维持其中少量假设。因而，模型将 ``任务中可能存在的规则'' 与 ``学习者当前正在考虑的规则'' 明确区分：$\mathcal H$ 表示长期可搜索空间，容量受限的活跃集合 $A_t\subset\mathcal H$ 表示试次 $t$ 实际参与推断的假设。

该模型由八个相互衔接的部分构成：
\begin{enumerate}
    \item 任务定义的几何假设空间；
    \item 基于独立知觉测量的被试特异知觉输入；
    \item 容量受限的活跃假设集合；
    \item 被试特异的 feedback-gated 多-profile假设转移；
    \item 由任务反馈事件定义的贝叶斯证据更新；
    \item 包含 fade-only 与 static-only 两个边界的双通道记忆；
    \item 假设特异的动态逆温度 $\beta_{t,h}$；
    \item 将内部假设信念映射为外显选择的连续 concentration readout。
\end{enumerate}
这些部分共同定义一个因果的状态滤波器。模型在作出试次 $t$ 的选择预测时只能读取试次 $t-1$ 及更早的反馈；当前反馈 $r_t$ 只能用于形成反馈后的状态，并预测下一试次。模型读入物理刺激后，先依据独立知觉任务得到的被试特异误差分布抽样知觉刺激；选择预测、反馈似然和内部学习均使用该知觉刺激。模型拟合还使用被试选择与任务反馈，口头规则报告不参与参数选择，而用于拟合后的外部效度检验。


\subsubsection{符号、观测与状态变量}

对被试 $s$ 的第 $t$ 个试次，观测记为
\begin{equation}
    o_{s,t}=(\mathbf x_{s,t},c_{s,t},r_{s,t}),
    \qquad
    \mathbf x_{s,t}\in[0,1]^4 ,
    \label{eq:modelv2-observation}
\end{equation}
其中 $\mathbf x_{s,t}$ 是归一化到 $[0,1]^4$ 的四维物理刺激，$c_{s,t}\in\{1,\ldots,N\}$ 是被试选择，$r_{s,t}$ 是任务反馈。知觉模块由 $\mathbf x_{s,t}$ 生成同一试次的知觉刺激 $\widetilde{\mathbf x}_{s,t}$。条件 1 中 $N=2$ 且 $r_t\in\{0,1\}$；条件 2 中 $N=4$ 且 $r_t\in\{0,1\}$；条件 3 中 $N=4$ 且 $r_t\in\{0,0.5,1\}$，分别表示类别家族错误、家族正确但精确类别错误，以及精确类别正确。

全文使用以下记号：
\begin{itemize}
    \item $\mathcal H=\{h_1,\ldots,h_K\}$：任务定义的完整候选假设空间；
    \item $\mathbf x_t$ 与 $\widetilde{\mathbf x}_t$：物理刺激及由被试特异知觉噪声生成的知觉刺激；
    \item $A_t\subset\mathcal H$：试次 $t$ 的活跃假设集合，满足 $|A_t|\leq M$；
    \item $m_t(h)=I(h\in A_t)$：活跃集合掩码；
    \item $\pi_t^-(h)$：试次 $t$ 反馈出现之前的假设先验；
    \item $\pi_t^+(h)$：整合试次 $t$ 反馈之后的假设后验；
    \item $q_{t,h}(i)=P(C_i\mid\widetilde{\mathbf x}_t,h,\beta_{t,h})$：假设 $h$ 对类别 $i$ 的预测概率；
    \item $\mathcal Y_t$：与试次 $t$ 的选择和反馈相容的真实类别集合；
    \item $L_t(h)=P(r_t\mid\widetilde{\mathbf x}_t,c_t,h,\beta_{t,h})$：反馈似然；
    \item $F_t(h)$ 与 $S_t(h)$：近期衰减证据和长期累积证据；
    \item $\beta_{t,h}$：假设 $h$ 在试次 $t$ 的逆温度；
    \item $\rho$：将假设先验读出为选择概率时的 concentration 参数。
\end{itemize}

$\pi_t^-$ 与 $\pi_t^+$ 必须在实现和结果文件中分别保存。$\pi_t^-$ 可以用于预测同一试次的选择；$\pi_t^+$ 已经包含当前反馈，只能解释反馈后的信念变化或预测下一试次。


\subsubsection{任务定义的几何假设空间}

\paragraph{统一表示。}
每个假设 $h\in\mathcal H$ 将四维单位超立方体 $[0,1]^4$ 划分为 $N$ 个带类别标签的区域。类别 $i$ 的区域写为
\begin{equation}
    R_{h,i}
    =
    \left\{
        \mathbf x\in[0,1]^4:
        A_{h,i}\mathbf x\leq\mathbf b_{h,i}
    \right\},
    \qquad i=1,\ldots,N .
    \label{eq:modelv2-region}
\end{equation}
每个区域对应一个代表性原型 $\boldsymbol\mu_{h,i}\in[0,1]^4$。区域定义保存规则的几何边界和类别标签；原型用于计算试次级类别概率。刺激按每名被试实际经历的 feature 顺序进入模型，因此假设的几何含义始终相对于该被试的四维坐标系定义。

\paragraph{条件 1：38 个带标签的二分类假设。}
与 V14 的实际配置一致，条件 1 首先生成 19 个基础几何划分：
\begin{enumerate}
    \item 单维轴向边界 $x_i=0.5$，共 4 个；
    \item 二维等值边界 $x_i=x_j$，共 $\binom{4}{2}=6$ 个；
    \item 二维求和边界 $x_i+x_j=1$，共 $\binom{4}{2}=6$ 个；
    \item 四维混合边界 $x_i+x_j=x_k+x_l$，共 3 个互不重复的维度配对。
\end{enumerate}
因此基础几何空间大小为
\begin{equation}
    |\mathcal H_1^{\mathrm{base}}|=4+6+6+3=19 .
    \label{eq:modelv2-cond1-count}
\end{equation}
V14 设置 \texttt{include\_label\_reversals=true}。对每个基础划分，模型同时保留恒等标签映射 $(0,1)$ 和反向标签映射 $(1,0)$；反向版本保持几何边界不变，但交换两个区域的类别标签及相应类别原型。因此，实际进入 V14 推断与搜索的完整带标签假设空间为
\begin{equation}
    \mathcal H_1
    =
    \mathcal H_1^{\mathrm{base}}
    \times
    \{(0,1),(1,0)\},
    \qquad
    |\mathcal H_1|=19\times2=38 .
    \label{eq:modelv2-cond1-labelled-count}
\end{equation}
轴向规则的两个类别原型在诊断维度上分别取 $0.25$ 和 $0.75$，其他维度取 $0.5$；等值规则使用 $(1/3,2/3)$ 与 $(2/3,1/3)$ 的镜像原型；求和规则在两个相关维度上共同取 $1/3$ 或 $2/3$；四维混合规则使等式两侧的维度采用相反的 $1/3$ 与 $2/3$ 排列。这里的 38 指带类别标签的假设，而不是 38 种不同的无标签几何边界。

\paragraph{条件 2 和 3：116 个四分类假设。}
条件 2 和 3 共享四分类几何空间。116 个候选划分的构成为
\begin{equation}
\begin{aligned}
    |\mathcal H_{2,3}|
    ={}&
    6_{\mathrm{axis+axis}}
    +6_{\mathrm{equality+sum}}
    +12_{\mathrm{axis+equality}}
    +12_{\mathrm{axis+sum}}\\
    &+3_{\mathrm{equality+equality}}
    +3_{\mathrm{sum+sum}}
    +24_{\mathrm{three\ axes}}\\
    &+24_{\mathrm{axis+equality+sum}}
    +12_{\mathrm{equality+two\ axes}}
    +12_{\mathrm{sum+two\ axes}}\\
    &+1_{\mathrm{dimension\ max}}
    +1_{\mathrm{dimension\ min}}
    =116 .
    \label{eq:modelv2-cond23-count}
\end{aligned}
\end{equation}
三轴划分以及轴向--等值--求和划分中超平面的顺序会改变四个区域的标签，因此在计数中保留顺序。dimension-max 假设把最大维度映射为类别，dimension-min 假设把最小维度映射为类别。

四分类原型与相应区域保持一致。轴向切分主要使用 $0.25/0.75$；等值与求和关系主要使用 $1/3$ 与 $2/3$；等值--求和产生的四个区域使用 $1/6$、$1/2$ 和 $5/6$；dimension-max 原型将目标维度设为 $0.8$、其余维度设为 $0.4$，dimension-min 原型采用相反设置。所有原型必须通过所属区域约束检查。

\paragraph{完整空间与有限活跃集合。}
$\mathcal H$ 只规定学习者可能搜索到哪些规则，并不表示这些规则同时处于工作状态。当前经验参照模型固定活跃容量 $M=5$，因此条件 1 每个试次最多维持 38 个带标签候选中的 5 个，条件 2 和 3 每个试次也只维持 116 个候选中的 5 个。扩大完整假设空间只增加长期搜索范围，不增加单个试次的并行假设负荷。有限容量是模型的理论承诺，但精确数值 $M=5$ 尚未通过容量恢复或相邻有限容量的公平比较独立确认；正式报告应将二者分开表述。

代码内部的整数索引只是枚举标签。保存或报告一个假设时，必须同时保存其边界、原型和类别映射，不能把 ``hypothesis 0'' 等实现索引当作可跨版本解释的理论名称。


\subsubsection{被试特异的知觉输入}

与 V14 的模块顺序一致，模型在假设转移、似然、记忆和 beta 更新之前首先运行知觉模块。该模块的参数不由当前类别学习行为拟合，而由每名被试在独立知觉任务中的误差估计给定；各维统计量按照该被试在 Task 2 中实际使用的四个 feature 顺序重新排列。

V14 的自动选择规则将被试 101--124、201--224 和 301--324 设为均匀阈限噪声模式。令 $\mathbf a_s=(a_{s,1},\ldots,a_{s,4})$ 为 \texttt{Task1b\_errorsummary\_72.csv} 中 \texttt{threshold\_mean\_mean} 的绝对值，则
\begin{align}
    \epsilon_{s,t,d}
    &\sim\operatorname{Uniform}(-a_{s,d},a_{s,d}),
    \qquad d=1,\ldots,4 .
    \label{eq:modelv2-perception-uniform}
\end{align}
被试 125--132、225--232、325--335 和 401--424 使用正态误差模式；其他未列入预设集合的被试也默认使用正态模式。令 $\mu_{s,d}$ 和 $\sigma_{s,d}$ 分别为 \texttt{Task1b\_errorsummary\_24.csv} 中按被试和 feature 汇总的 \texttt{error\_mean} 与 \texttt{error\_std}，则
\begin{equation}
    \epsilon_{s,t,d}
    \sim
    \mathcal N(\mu_{s,d},\sigma_{s,d}^{2}) .
    \label{eq:modelv2-perception-normal}
\end{equation}
每次模拟重复、每个试次均重新抽样噪声，并生成
\begin{equation}
    \widetilde{\mathbf x}_{s,t}
    =
    \operatorname{clip}
    \left(
        \mathbf x_{s,t}+\boldsymbol\epsilon_{s,t},
        0,1
    \right),
    \label{eq:modelv2-perceived-stimulus}
\end{equation}
其中截断逐维执行。因而，$\mathbf x_{s,t}$ 是实验呈现的物理刺激，$\widetilde{\mathbf x}_{s,t}$ 才是进入当前模拟轨迹的内部知觉输入。知觉抽样是模型随机性的一部分，最终边际选择概率同时积分知觉噪声、active-set 初始化、profile 抽样和 newcomer 抽样。


\subsubsection{单个假设下的类别概率}

对假设 $h$ 和类别 $i$，知觉刺激到类别原型的欧氏距离为
\begin{equation}
    d_{t,h,i}
    =
    \left\|
        \widetilde{\mathbf x}_t-\boldsymbol\mu_{h,i}
    \right\|_2 .
    \label{eq:modelv2-prototype-distance}
\end{equation}
假设 $h$ 通过距离 softmax 产生类别概率：
\begin{equation}
    q_{t,h}(i)
    =
    \frac{
        \exp[-\beta_{t,h}d_{t,h,i}]
    }{
        \sum_{j=1}^{N}
        \exp[-\beta_{t,h}d_{t,h,j}]
    } .
    \label{eq:modelv2-category-softmax}
\end{equation}
$\beta_{t,h}$ 越大，类别概率越集中于最近原型；$\beta_{t,h}$ 越小，概率越接近均匀分布。$\beta_{t,h}$ 控制单个假设内部的类别边界锐度，而 $\rho$ 控制多个假设之间的 readout 集中程度，两者承担不同作用。


\subsubsection{逐试次因果顺序}

模型在每个试次按照以下顺序运行：
\begin{enumerate}
    \item 读取物理刺激 $\mathbf x_t$，按被试特异误差分布抽样并截断得到 $\widetilde{\mathbf x}_t$；
    \item 读取上一试次的活跃集合 $A_{t-1}$、后验 $\pi_{t-1}^+$、记忆状态和 $\beta_{t-1,h}$；
    \item 只使用截至 $t-1$ 的反馈历史和 $\pi_{t-1}^+$，抽取当前 profile $K_t$；
    \item 根据 $K_t$ 生成容量不超过 $M$ 的 $A_t$，并将 $\pi_{t-1}^+$ 映射为 $\pi_t^-$；
    \item 使用 $(\widetilde{\mathbf x}_t,\pi_t^-,\beta_{t,h})$ 计算当前选择概率；
    \item 在观察真实选择 $c_t$ 和反馈 $r_t$ 后，构造反馈事件 $\mathcal Y_t$ 和似然 $L_t(h)$；
    \item 将 $L_t(h)$ 写入双通道记忆，得到 $\pi_t^+$；
    \item 先保存本试次实际使用的 $\beta_{t,h}$，再根据当前反馈更新 $\beta_{t+1,h}$；
    \item 将 $r_t$ 写入历史，供试次 $t+1$ 的 profile controller 使用。
\end{enumerate}

因此，试次 $t$ 的选择预测满足
\begin{equation}
    P(c_t=i\mid\mathcal D_{1:t-1},\widetilde{\mathbf x}_t)
    =
    \sum_{h\in A_t}
    \omega_{t,h}(\rho)
    q_{t,h}(i),
    \label{eq:modelv2-causal-choice}
\end{equation}
其中 $\mathcal D_{1:t-1}$ 只包含过去试次。当前反馈不能进入 $\pi_t^-$、profile $K_t$ 或产生当前选择的 $\beta_{t,h}$。

该架构是以真实行为序列为条件的一步状态滤波器：模型先预测当前选择，然后使用观察到的选择和反馈更新内部状态。若要生成完全新的自主行为轨迹，则需在该滤波器外部由式 \ref{eq:modelv2-causal-choice} 抽取选择、由实验任务规则生成反馈，再调用相同的状态更新。


\subsubsection{被试特异的多-profile有限容量假设转移}

\paragraph{初始化。}
第一个试次从 $\mathcal H$ 中无放回抽取 $M=5$ 个假设：
\begin{equation}
    A_1\sim\operatorname{UniformSubset}(\mathcal H,M),
    \qquad
    \pi_1^-(h)=
    \begin{cases}
        1/M,&h\in A_1,\\
        0,&h\notin A_1.
    \end{cases}
    \label{eq:modelv2-initial-active-set}
\end{equation}
初始化随机性通过重复模拟积分，不能以单一有利初始集合代表一个参数配置。

\paragraph{控制信号。}
单一、固定的 profile 无法表达被试在规则维持、局部探索和较大幅度修正之间的差异。模型因此保留一个被试特异的 feedback-gated controller。Controller 只读取过去反馈、上一后验和试次进程，不读取当前反馈。令
\begin{align}
    e_{t-1}&=1-r_{t-1},\\
    \bar a_t
    &=
    \frac{1}{W}
    \sum_{j=t-W}^{t-1}\tilde r_j,
    \qquad
    \bar e_t=1-\bar a_t,\\
    \tau_t
    &=
    \min\left(\frac{t-1}{160},1\right),
    \label{eq:modelv2-controller-history}
\end{align}
其中 $W=8$；条件 3 的部分反馈 $r_j=0.5$ 对 $\bar a_t$ 贡献 $0.5$。历史不足 $W$ 个试次时，用该实验条件下随机选择的期望反馈填充缺失的 $r_j$，而不是使用当前反馈。
具体地，条件 1 和条件 2 的填充值为
\[
    r_{\mathrm{chance}}=\frac{1}{N};
\]
若条件 3 的每个类别家族包含两个类别，则随机选择获得精确反馈的概率为 $1/N$，获得家族反馈的概率也为 $1/N$，所以
\[
    r_{\mathrm{chance}}
    =
    \frac{1}{N}+\frac{0.5}{N}
    =
    \frac{3}{2N}
    =0.375 .
\]

上一后验的归一化熵为
\begin{equation}
    H_t
    =
    \begin{cases}
    \displaystyle
    \frac{
        -\sum_{h\in A_{t-1}}
        \pi_{t-1}^+(h)\log\pi_{t-1}^+(h)
    }{
        \log |A_{t-1}|
    },
    &|A_{t-1}|>1,\\[12pt]
    0,&|A_{t-1}|=1.
    \end{cases}
    \label{eq:modelv2-active-entropy}
\end{equation}
熵按实际活跃集合大小归一化，因此 $H_t\in[0,1]$：$H_t$ 高表示模型在当前活跃假设之间不确定，$1-H_t$ 高表示后验集中。

\paragraph{Feedback-gated profile 概率。}
当前经验参照模型包含四种操作 profile：
\[
    K_t\in\{\mathrm C,\mathrm S,\mathrm A,\mathrm B\},
\]
分别表示 Conservative、Stable、Aggressive 和 Stubborn；下标 $\mathrm B$ 用于避免与 Stable 的 $\mathrm S$ 混淆。这些名称表示算法行为，不表示固定的人格类型。每名被试选择一个 controller family $g_s$。在 family $g_s$ 下，profile $k$ 的激活值和抽样概率为
\begin{equation}
    u_{t,k}^{(g_s)}
    =
    \mathbf a_k^{(g_s)\top}\mathbf f_t,
    \qquad
    P(K_t=k\mid g_s)
    =
    \frac{
        \exp[u_{t,k}^{(g_s)}/T_{g_s}]
    }{
        \sum_{\ell\in\{\mathrm C,\mathrm S,\mathrm A,\mathrm B\}}
        \exp[u_{t,\ell}^{(g_s)}/T_{g_s}]
    },
    \label{eq:modelv2-profile-softmax}
\end{equation}
其中
\[
    \mathbf f_t
    =
    (1,\bar a_t,\bar e_t,e_{t-1},H_t,1-H_t,\tau_t)
\]
且 $T_{g_s}>0$。因而，同一被试可以在不同试次进入不同 profile；subject-specific 指的是 controller 的激活结构，而不是把被试永久归入某一种 profile。

\paragraph{六个 controller family。}
当前结构筛选保留六个具有不同反馈反应模式的 family：
\begin{enumerate}
    \item \texttt{stable\_dominant}：以 Stable 为主要状态，同时允许错误和高熵提高 Aggressive；
    \item \texttt{error\_aggressive}：错误、近期低正确率和高熵更强地提高 Aggressive；
    \item \texttt{early\_explore\_late\_stable}：学习早期更易探索，后期更易稳定维持；
    \item \texttt{conservative\_heavy}：高正确率和高置信度更强地维持现有集合；
    \item \texttt{error\_choice\_newcomer}：错误后提高 Stubborn，并允许近期错误选择参与候选打分；
    \item \texttt{choice\_volatile\_refresh}：保留 choice-informed 打分，同时提高 profile 抽样和集合刷新的变异性。
\end{enumerate}
这六个 family 是被试层面的离散候选，不是六个同时叠加的模块。令
\[
    A=\bar a_t,\quad E=\bar e_t,\quad
    e=e_{t-1},\quad H=H_t,\quad
    C=1-H_t,\quad \tau=\tau_t .
\]
六个 family 对 $(\mathrm C,\mathrm S,\mathrm A,\mathrm B)$ 的当前经验参照 logits 为
\begin{align}
    \texttt{stable\_dominant}:\quad
    (u_C,u_S,u_A,u_B)
    ={}&
    (0.1+0.7A+0.5C,\;
     1.0+0.5H+0.3A,\nonumber\\
    &-0.4+0.7E+0.6H,\;
     -0.3+0.7e+0.5C),\\
    \texttt{error\_aggressive}:\quad
    (u_C,u_S,u_A,u_B)
    ={}&
    (-0.2+0.9A+0.6C,\;
     0.2+0.4H+0.4A,\nonumber\\
    &0.3+1.2e+1.2E+0.8H,\;
     -0.5+0.3e+0.4C),\\
    \texttt{early\_explore\_late\_stable}:\quad
    (u_C,u_S,u_A,u_B)
    ={}&
    (-0.2+1.2\tau+0.8A,\;
     0.3+0.8\tau+0.5H,\nonumber\\
    &0.5+H-1.2\tau,\;
     -0.3+0.8e+0.5C),\\
    \texttt{conservative\_heavy}:\quad
    (u_C,u_S,u_A,u_B)
    ={}&
    (1.0+A+C,\;
     0.2+0.6H,\;
     -1.0+0.5E,\;
     -0.6+0.5e),\\
    \texttt{error\_choice\_newcomer}:\quad
    (u_C,u_S,u_A,u_B)
    ={}&
    (-1.5+A-0.5E+0.8C,\;
     -1.0+0.8A+0.4C-0.5H,\nonumber\\
    &-1.2+0.5e+E+H-0.5C,\;
     0.8+0.7e+0.4E+C),\\
    \texttt{choice\_volatile\_refresh}:\quad
    (u_C,u_S,u_A,u_B)
    ={}&
    (-1.5+A-0.5E+0.8C,\;
     -0.2+0.9H+0.2A,\nonumber\\
    &0.9e+1.4E+H,\;
     -0.1+0.8e+0.4E+0.6C).
    \label{eq:modelv2-controller-family-logits}
\end{align}
相应温度依次为 $0.85$、$0.70$、$0.80$、$0.75$、$0.65$ 和 $1.10$。由于式 \ref{eq:modelv2-active-entropy} 使用 active-set entropy，而旧实现曾按完整假设空间归一化熵，上述系数在正式 v2 中必须作为重新优化的起点，而不能未经重拟合直接冻结。

\paragraph{四种操作 profile。}
令上一活跃假设按其 survivor 得分排序。四种 profile 对集合的操作为：
\begin{enumerate}
    \item \textit{Conservative}：保持 $A_{t-1}$ 不变，不引入 newcomer；
    \item \textit{Stable}：通常保留 $M-1$ 个旧假设并引入 1 个 newcomer；旧假设按 survivor 得分经温度变换后无放回抽样；
    \item \textit{Aggressive}：保留得分最高的一个旧假设 $h^\star$，并引入
    \begin{equation}
        n_t^{\mathrm{new}}
        =
        \operatorname{clip}
        \left[
            \operatorname{round}
            \{(1-\pi_{t-1}^+(h^\star))n_{\max}\},
            n_{\min},n_{\max}
        \right]
        \label{eq:modelv2-aggressive-count}
    \end{equation}
    个 newcomer；
    \item \textit{Stubborn}：确定性保留得分最高的 $n_{\mathrm{retain}}$ 个旧假设，并以
    \begin{equation}
        p_t^{\mathrm{explore}}
        =
        p_0+p_{\mathrm{correct}}(1-e_{t-1})
        +p_{\mathrm{error}}e_{t-1}
        \label{eq:modelv2-stubborn-explore}
    \end{equation}
    的概率引入至多一个 newcomer。
\end{enumerate}
任一 profile 都至少保留一个旧假设；若 inactive pool 不足，则使用全部剩余候选并保持 $|A_t|\leq M$。当前六个 family 均使用 $M=5$；Stable 的 newcomer 数为 1，Aggressive 的 $n_{\max}=4$，Stubborn 的 $n_{\mathrm{retain}}=2$。不同 family 只改变保留温度、最小 newcomer 数、探索概率和 newcomer 质量，而不改变四种 profile 的基本语义。

定义幸存者和新进入者为
\begin{equation}
    \mathcal S_t=A_{t-1}\cap A_t,
    \qquad
    \mathcal N_t=A_t\setminus A_{t-1},
    \qquad
    n_t^{\mathrm{new}}=|\mathcal N_t|.
    \label{eq:modelv2-survivor-newcomer}
\end{equation}

\paragraph{Survivor 与 newcomer 得分。}
默认 survivor 得分为上一后验：
\[
    s_t^{\mathrm{surv}}(h)=\pi_{t-1}^+(h).
\]
\texttt{error\_choice\_newcomer} 和 \texttt{choice\_volatile\_refresh} 允许上一选择的可解释性参与 survivor 得分：
\begin{equation}
    s_t^{\mathrm{surv}}(h)
    \propto
    [\pi_{t-1}^+(h)]^a
    [q_{t-1,h}(c_{t-1})]^b .
    \label{eq:modelv2-choice-survivor}
\end{equation}
对未激活假设，这两个 family 使用最近 $W_c=8$ 个错误试次的选择概率：
\begin{equation}
    \log s_t^{\mathrm{new}}(h)
    =
    \frac{b_c}{|\mathcal E_t|}
    \sum_{j\in\mathcal E_t}
    \log q_{j,h}(c_j),
    \qquad
    \mathcal E_t=
    \{j<t:r_j<1\}\cap\{t-W_c,\ldots,t-1\}.
    \label{eq:modelv2-choice-newcomer}
\end{equation}
其他 family 从 inactive pool 均匀抽取 newcomer。Choice-informed 得分只属于相应 controller family，不是所有被试共享的默认状态变量。

\paragraph{转移后的先验质量。}
令 $\nu_{t,k}\in[0,1]$ 为 profile $k$ 分给 newcomer 的总质量。给定 profile、survivor 和 newcomer 得分，转移先验统一写为
\begin{equation}
    \pi_t^-(h)
    =
    \begin{cases}
    \displaystyle
    (1-\nu_{t,k})
    \frac{s_t^{\mathrm{surv}}(h)}
    {\sum_{u\in\mathcal S_t}s_t^{\mathrm{surv}}(u)},
    &h\in\mathcal S_t,\\[14pt]
    \displaystyle
    \nu_{t,k}
    \frac{s_t^{\mathrm{new}}(h)}
    {\sum_{u\in\mathcal N_t}s_t^{\mathrm{new}}(u)},
    &h\in\mathcal N_t,\\[12pt]
    0,&h\notin A_t .
    \end{cases}
    \label{eq:modelv2-transition-prior}
\end{equation}
Conservative 使用 $\nu_{t,C}=0$；Aggressive 将 $1-\pi_{t-1}^+(h^\star)$ 分给 newcomer；Stubborn 使用 family-specific 的小质量。Stable 使用相似性--新颖性映射：高 posterior confidence 时优先给与旧解释相似的 newcomer 质量，低 confidence 时提高结构上新颖候选的质量。假设相似性通过在 $[0,1]^4$ 上预先抽取 $B=100{,}000$ 个共同点并比较带标签类别分配计算：
\begin{equation}
    \operatorname{Sim}(h,u)
    =
    \frac{1}{B}
    \sum_{b=1}^{B}
    I\{g_h(\mathbf x_b)=g_u(\mathbf x_b)\}.
    \label{eq:modelv2-hypothesis-similarity}
\end{equation}

\paragraph{证据边界与精简条件。}
现有配对删除结果支持多-profile controller 作为整体保留，但不支持把每个内部 profile 都解释为同等必要。删除 Stable 的平均 Brier 代价为 $0.0203$；Stubborn 和 Conservative 分别为 $0.0048$ 和 $0.0031$；当前 Aggressive 的平均贡献为 $-0.0047$。另一方面，把 subject-specific controller 统一为 stable-dominant 在既有结构筛选中造成约 $0.0147$--$0.0183$ 的平均 Brier 损失。因此，正式 v2 在完成直接比较之前保留六-family、四-profile controller 作为经验参照，不能用未经拟合的
\[
    P(S)=1-\bar e_t,\qquad
    P(P)=\bar e_t(1-H_t),\qquad
    P(R)=\bar e_tH_t
\]
替代它。三-profile 压缩只有在与六-family参照使用相同预算、联合重优化其他参数后，在边际 Brier、轨迹 CRPS 和被试级预测上达到预设非劣标准，才能取代当前 controller。代码重构可以合并重复实现，但不能在这一比较之前删除相应的行为自由度。


\subsubsection{反馈事件与贝叶斯似然}

\paragraph{反馈相容类别集合。}
模型首先把任务反馈转换为一组互斥、完备的反馈事件。条件 1 和条件 2 中，
\begin{equation}
    \mathcal Y_t
    =
    \begin{cases}
        \{c_t\},&r_t=1,\\
        \{1,\ldots,N\}\setminus\{c_t\},&r_t=0.
    \end{cases}
    \label{eq:modelv2-feedback-set-cond12}
\end{equation}

条件 3 使用由实验设计固定、对所有假设相同的类别家族映射 $\mathcal G(c)$。若类别编码中的两个家族为 $\{1,2\}$ 和 $\{3,4\}$，则 $\mathcal G(1)=\mathcal G(2)=\{1,2\}$，$\mathcal G(3)=\mathcal G(4)=\{3,4\}$；若实际数据采用其他标签顺序，应在读入数据时显式重映射，而不能让家族关系随假设几何改变。反馈事件为
\begin{equation}
    \mathcal Y_t
    =
    \begin{cases}
        \{c_t\},&r_t=1,\\
        \mathcal G(c_t)\setminus\{c_t\},&r_t=0.5,\\
        \{1,\ldots,4\}\setminus\mathcal G(c_t),&r_t=0.
    \end{cases}
    \label{eq:modelv2-feedback-set-cond3}
\end{equation}
三个集合互不重叠，并完整覆盖所有真实类别，因此 $r_t=0.5$ 与 $r_t=0$ 表示不同的实验事件。

\paragraph{反馈似然。}
假设 $h$ 对观测反馈的支持为
\begin{equation}
    L_t(h)
    =
    \sum_{i\in\mathcal Y_t}q_{t,h}(i).
    \label{eq:modelv2-feedback-likelihood}
\end{equation}
计算时将 $L_t(h)$ 截断到 $[\epsilon,1-\epsilon]$ 以避免对数溢出，但截断常数 $\epsilon$ 必须固定并报告。对条件 1，错误反馈的 $\mathcal Y_t$ 只有另一个类别；对条件 2，错误反馈的似然是其余三个类别概率之和；对条件 3，精确正确、家族正确和家族错误分别累加三个互斥集合中的概率。

若不包含记忆扩展，一步规范性更新为
\begin{equation}
    \pi_t^+(h)
    =
    \frac{
        \pi_t^-(h)L_t(h)m_t(h)
    }{
        \sum_{u\in\mathcal H}
        \pi_t^-(u)L_t(u)m_t(u)
    }.
    \label{eq:modelv2-standard-bayes}
\end{equation}
正式模型通过下一节的双通道证据状态实现跨试次积累，但仍保持 ``先验由转移产生、当前反馈通过似然进入后验'' 的贝叶斯骨架。


\subsubsection{包含 static-only 边界的双通道记忆}

\paragraph{两条证据轨道。}
对每个活跃假设，模型维护近期衰减轨道 $F_t(h)$ 和长期累积轨道 $S_t(h)$：
\begin{align}
    F_t(h)
    &=\gamma F_{t^-}(h)+\log L_t(h),\\
    S_t(h)
    &=S_{t^-}(h)+\log L_t(h),
    \label{eq:modelv2-dual-memory}
\end{align}
其中 $t^-$ 表示当前反馈进入之前、完成 active-set 转移和状态同步之后的状态，$0\leq\gamma<1$。两条轨道按
\begin{equation}
    \ell_t(h)
    =
    w_0S_t(h)+(1-w_0)F_t(h),
    \qquad 0\leq w_0\leq1 ,
    \label{eq:modelv2-memory-mixture}
\end{equation}
混合，并在 $A_t$ 内归一化：
\begin{equation}
    \pi_t^+(h)
    =
    \frac{
        \exp[\ell_t(h)]m_t(h)
    }{
        \sum_{u\in\mathcal H}
        \exp[\ell_t(u)]m_t(u)
    } .
    \label{eq:modelv2-memory-posterior}
\end{equation}
第一个试次在当前活跃集合上初始化为
\begin{equation}
    F_{1^-}(h)=S_{1^-}(h)=\log\pi_1^-(h),
    \qquad h\in A_1.
    \label{eq:modelv2-memory-initialization}
\end{equation}

$w_0=0$ 时，式 \ref{eq:modelv2-memory-mixture} 只读取衰减轨道，模型退化为 fade-only；$w_0=1$ 时只读取累积轨道，模型严格退化为 static-only；只有 $0<w_0<1$ 时两条轨道同时影响后验。因此，fade-only、static-only 和内部混合不是三个互不相关的架构，而是同一个双通道模型族中的两个边界和内部区域。

\paragraph{转移后的状态同步。}
Hypothesis transition 已经通过式 \ref{eq:modelv2-transition-prior} 定义新的 $\pi_t^-$，所以双通道内部状态必须在吸收当前似然之前与该先验一致。对幸存假设，保留其长期与近期轨道差
\begin{equation}
    \Delta_{t,h}
    =
    S_{t-1}(h)-F_{t-1}(h).
    \label{eq:modelv2-memory-offset}
\end{equation}
令
\begin{equation}
    T_{t,h}=\log\pi_t^-(h)+C_t ,
    \label{eq:modelv2-memory-target}
\end{equation}
其中 $C_t$ 是对所有活跃假设相同、归一化后抵消的平移常数。同步状态定义为
\begin{align}
    F_{t^-}(h)
    &=T_{t,h}-w_0\Delta_{t,h},\\
    S_{t^-}(h)
    &=T_{t,h}+(1-w_0)\Delta_{t,h}.
    \label{eq:modelv2-memory-sync}
\end{align}
因此
\[
    w_0S_{t^-}(h)+(1-w_0)F_{t^-}(h)=T_{t,h},
\]
同时保留幸存假设的长短期证据差。对 newcomer，设 $\Delta_{t,h}=0$，即
\[
    F_{t^-}(h)=S_{t^-}(h)=T_{t,h};
\]
对离开 active set 的假设，两条轨道设为 $-\infty$。同步必须在每个试次、当前似然进入之前执行，而不只在 active-set 掩码发生变化时执行。

\paragraph{参数边界与可识别性。}
双通道拟合必须显式允许
\begin{equation}
    w_0\in[0,1]
    \label{eq:modelv2-w0-range}
\end{equation}
并实际评估 $w_0=0$ 与 $w_0=1$。遗漏 $w_0=1$ 会使被拟合的 ``双通道模型'' 不再包含 static-only，从而破坏嵌套模型比较。当 $w_0=1$ 时，$\gamma$ 不影响后验，应记录为不可识别或不适用，而不能解释其估计值。当 $\gamma=1$ 且两条轨道初始化相同时，两条轨道完全重合，$w_0$ 同样不可识别；因此内部双通道解要求 $0<w_0<1$ 且 $\gamma<1$。

在相同训练损失、参数空间和充分优化条件下，包含 $w_0=1$ 的双通道模型族的最优训练损失不可能高于 static-only。是否需要内部双通道解，应根据重新优化后的参数位置、独立预测表现和复杂度惩罚共同判断，而不能用冻结其他参数后的一次局部切换替代完整比较。


\subsubsection{假设特异的动态 beta}

\paragraph{初始化与解释。}
每个活跃假设拥有独立的 $\beta_{t,h}$。初始假设和 newcomer 均使用
\begin{equation}
    \beta_{t,h}=\beta_0,
    \qquad
    \beta_{\min}\leq\beta_0\leq\beta_{\max}.
    \label{eq:modelv2-beta-init}
\end{equation}
当前固定值为
\[
    \beta_0=5,\qquad
    \beta_{\min}=0.1,\qquad
    \beta_{\max}=25.
\]
假设离开 active set 后不继续更新；若以后重新进入，则重新设为 $\beta_0$。这样，$\beta_{t,h}$ 表示学习者当前使用假设 $h$ 时，其刺激到类别映射的确定程度，而不是该假设的 posterior mass。

\paragraph{条件通用的反馈一致性。}
为使动态 beta 在二分类、四分类和家族反馈条件下具有相同含义，首先使用式 \ref{eq:modelv2-feedback-likelihood} 定义假设对观测反馈事件的支持：
\[
    s_{t,h}=L_t(h).
\]
在均匀类别预测下，该事件的机会概率为
\begin{equation}
    b_t=\frac{|\mathcal Y_t|}{N}.
    \label{eq:modelv2-feedback-chance}
\end{equation}
将支持度相对于机会水平标准化为 $[-1,1]$：
\begin{equation}
    \xi_{t,h}
    =
    \begin{cases}
    \displaystyle
    \frac{s_{t,h}-b_t}{1-b_t},
    &s_{t,h}\geq b_t,\\[10pt]
    \displaystyle
    \frac{s_{t,h}-b_t}{b_t},
    &s_{t,h}<b_t.
    \end{cases}
    \label{eq:modelv2-beta-evidence}
\end{equation}
$\xi_{t,h}>0$ 表示假设给观测反馈事件分配了高于机会水平的概率，$\xi_{t,h}<0$ 表示低于机会水平。

Beta 在当前反馈之后更新：
\begin{equation}
\begin{aligned}
    \beta_{t+1,h}
    =
    \operatorname{clip}\Bigg[
    &\beta_{t,h}
    +\eta_+\max(\xi_{t,h},0)
    \frac{\beta_{\max}-\beta_{t,h}}{\beta_{\max}}\\
    &-\eta_-\max(-\xi_{t,h},0)\beta_{t,h},
    \ \beta_{\min},\beta_{\max}
    \Bigg].
    \label{eq:modelv2-beta-update}
\end{aligned}
\end{equation}
当前学习率为 $\eta_+=0.5$ 和 $\eta_-=0.15$。正证据使类别映射逐渐锐化，且接近上界时增量减小；负证据使映射按当前 beta 的比例变平。由于 $\mathcal Y_t$ 已经编码条件 3 的固定家族结构，$r_t=0.5$ 可以直接改变 beta，不需要从错误反馈猜测一个唯一正确类别。

\paragraph{因果记录。}
式 \ref{eq:modelv2-beta-update} 必须在使用 $\beta_{t,h}$ 计算当前 $q_{t,h}$、选择概率和 $L_t(h)$ 之后执行。结果文件分别保存 pre-feedback $\beta_{t,h}$ 和更新后的 $\beta_{t+1,h}$；任何同试次选择分析都只能使用前者。

\paragraph{证据边界。}
现有消融支持保留 ``反馈驱动、假设特异的动态 beta'' 这一机制：它相对固定 $\beta=5$ 的平均 Brier 收益为 $0.0542$，8/8 代表性被试方向一致。然而，这一结果尚未排除被试特异的最优静态 beta，且式 \ref{eq:modelv2-beta-evidence}--\ref{eq:modelv2-beta-update} 的条件通用版本需要重新拟合。因此，主模型保留动态 beta 的功能自由度，但不把当前学习率或新多分类更新式描述为已经由同一消融单独确认。


\subsubsection{连续 concentration readout}

内部先验 $\pi_t^-$ 表示当前选择之前对假设的信念，但外显选择未必对所有活跃假设进行线性平均。模型用一个连续 concentration 参数 $\rho\geq1$ 定义 readout 权重：
\begin{equation}
    \omega_{t,h}(\rho)
    =
    \frac{
        [\pi_t^-(h)]^\rho
    }{
        \sum_{u\in A_t}[\pi_t^-(u)]^\rho
    },
    \qquad h\in A_t .
    \label{eq:modelv2-readout-weight}
\end{equation}
最终类别概率为
\begin{equation}
    \widehat P_t(C_i)
    =
    \sum_{h\in A_t}
    \omega_{t,h}(\rho)q_{t,h}(i).
    \label{eq:modelv2-choice-probability}
\end{equation}
$\rho=1$ 是 posterior expectation；$\rho>1$ 逐渐增加高概率假设对选择的支配程度；当 $\rho\rightarrow\infty$ 时，readout 接近只使用最大 posterior 假设。Expectation、sharpened expectation 和 MAP 极限因此属于同一连续机制，而不是相互独立的输出模块。

Readout 只改变由 $\pi_t^-$ 到选择概率的映射，不回写 active set、profile、likelihood、memory 或 beta。在其他参数和随机种子相同的条件下，不同 $\rho$ 共享完全相同的潜在学习状态轨迹。拟合时必须把 $\rho=1$ 作为显式边界，并预先报告连续搜索采用的数值上界或变换方式。

现有消融支持 readout 灵活性：相对 expectation，所选灵活 readout 的平均 Brier 收益为 $0.0233$，8/8 代表性被试方向一致。该证据来自离散 expectation、sharpened 和 MAP 候选；连续 $\rho$ 是对这些候选的统一参数化，仍需在联合重优化中确认，而不能把离散候选的收益直接当作连续参数估计的验证。


\subsubsection{完整逐试次算法}

对一次模拟重复，完整算法如下：
\begin{enumerate}
    \item 从 $\mathcal H$ 中抽取 $A_1$，令 $\pi_1^-$ 在 $A_1$ 上均匀，并令所有初始假设的 $\beta_{1,h}=\beta_0$；
    \item 对试次 $t=1,\ldots,T$：
    \begin{enumerate}
        \item 由物理刺激 $\mathbf x_t$ 按式 \ref{eq:modelv2-perception-uniform} 或 \ref{eq:modelv2-perception-normal} 抽样噪声，并按式 \ref{eq:modelv2-perceived-stimulus} 得到 $\widetilde{\mathbf x}_t$；
        \item 若 $t=1$，直接使用初始化得到的 $A_1$ 和 $\pi_1^-$；若 $t>1$，由过去反馈和 $\pi_{t-1}^+$ 构造 controller features，按式 \ref{eq:modelv2-profile-softmax} 在被试 family $g_s$ 下抽取 $K_t$，再由对应 profile 和式 \ref{eq:modelv2-transition-prior} 生成 $A_t$ 与 $\pi_t^-$；
        \item 按式 \ref{eq:modelv2-category-softmax} 计算每个活跃假设的类别概率；
        \item 按式 \ref{eq:modelv2-readout-weight}--\ref{eq:modelv2-choice-probability} 计算当前选择预测；
        \item 读取真实 $c_t,r_t$，由式 \ref{eq:modelv2-feedback-set-cond12} 或 \ref{eq:modelv2-feedback-set-cond3} 构造 $\mathcal Y_t$；
        \item 按式 \ref{eq:modelv2-feedback-likelihood} 计算反馈似然；
        \item 同步两条记忆轨道，再按式 \ref{eq:modelv2-dual-memory}--\ref{eq:modelv2-memory-posterior} 得到 $\pi_t^+$；
        \item 保存 pre-feedback $\beta_{t,h}$，再按式 \ref{eq:modelv2-beta-update} 得到 $\beta_{t+1,h}$；
        \item 保存 $A_t,K_t,\pi_t^-,\pi_t^+,F_t,S_t,\beta_t$ 和选择预测。
    \end{enumerate}
\end{enumerate}
所有概率归一化都只在当前 active set 内进行。若数值截断导致总质量为 0，则使用当前 active set 上的均匀分布作为明确记录的数值兜底；该事件的发生次数必须进入诊断日志。


\subsubsection{被试层参数、随机性与拟合目标}

\paragraph{参数层级。}
精简模型固定任务结构、$M=5$、$W=8$ 以及 dynamic-beta 更新常数。主要被试层参数为
\begin{equation}
    \Omega_s=(g_s,\gamma_s,w_{0,s},\rho_s).
    \label{eq:modelv2-subject-parameters}
\end{equation}
其中 $g_s$ 是六个 controller family 之一，$\gamma_s$ 控制近期证据衰减，$w_{0,s}$ 控制长期与近期证据的混合，$\rho_s$ 控制外显选择对高 posterior 假设的集中程度。Controller family 必须与 memory 和 readout 联合重优化；不能先在一组固定参数下选择 $g_s$，再把由此产生的局部最优当作其他模块的独立效应。

\paragraph{随机变量与边际预测。}
知觉噪声、初始化、profile 抽样和 newcomer 抽样使同一参数组合产生不同潜在轨迹。对每个参数组合运行 $R$ 次模拟，正式评分使用边际类别概率
\begin{equation}
    \bar P_t(C_i\mid\Omega_s)
    =
    \frac{1}{R}
    \sum_{r=1}^{R}
    \widehat P_{t,r}(C_i\mid\Omega_s).
    \label{eq:modelv2-marginal-choice}
\end{equation}
参数比较应使用相同随机数流，使两个候选之间的差异尽可能来自参数或结构改变，而不是不同的随机初始集合。搜索阶段与最终确认阶段分别报告 $R$；最终 $R$ 应通过边际预测和损失的 Monte Carlo 收敛检查确定。

\paragraph{Choice Brier 与 NLL。}
主拟合损失为多类别 Brier score：
\begin{equation}
    \mathcal L_{\mathrm{Brier}}(\Omega_s)
    =
    \frac{1}{|\mathcal T_s|}
    \sum_{t\in\mathcal T_s}
    \sum_{i=1}^{N}
    \left[
        \bar P_t(C_i\mid\Omega_s)-I(c_t=i)
    \right]^2 .
    \label{eq:modelv2-brier}
\end{equation}
同时报告选择负对数似然
\begin{equation}
    \mathcal L_{\mathrm{NLL}}(\Omega_s)
    =
    -\frac{1}{|\mathcal T_s|}
    \sum_{t\in\mathcal T_s}
    \log\bar P_t(C_{c_t}\mid\Omega_s).
    \label{eq:modelv2-nll}
\end{equation}
第一个试次主要用于随机状态初始化，是否计入 $\mathcal T_s$ 必须在所有模型和被试间保持一致。

\paragraph{轨迹校准。}
为避免只拟合平均选择概率，模型还评价滑动窗口正确率的经验 continuous ranked probability score（CRPS）。设第 $j$ 个窗口的真实正确率为 $y_j$，第 $r$ 次模拟的预测正确率为 $\hat y_{j,r}$，则
\begin{equation}
    \operatorname{CRPS}_j
    =
    \frac{1}{R}
    \sum_{r=1}^{R}|\hat y_{j,r}-y_j|
    -
    \frac{1}{2R^2}
    \sum_{r=1}^{R}\sum_{r'=1}^{R}
    |\hat y_{j,r}-\hat y_{j,r'}|.
    \label{eq:modelv2-crps}
\end{equation}
最终 CRPS 对窗口求平均。Brier 决定主要参数选择，NLL 检查概率分配是否出现灾难性低值，CRPS 检查预测轨迹的位置和离散度。三项指标必须基于全部模拟的边际或经验分布，而不是选择表现最好的若干随机轨迹。

\paragraph{双通道的公平搜索。}
$w_0$ 的候选或连续优化区间必须包含 0 和 1。对 $w_0=1$ 的候选不搜索或解释 $\gamma$；对内部解 $0<w_0<1$，联合重新优化 $\gamma$、$w_0$ 和 $\rho$。所有候选使用相同模拟预算和随机数流。这样，static-only 是双通道模型族内的合法边界，而不是使用不同优化预算的外部模型。


\subsubsection{口头规则报告的独立验证}

口头报告不参与 $\Omega_s$ 的估计。若报告在被试作出选择之后、反馈出现之前获得，则模型比较状态应吸收被试已经作出的选择，但不能吸收当前反馈：
\begin{equation}
    \pi_t^{\mathrm{oral}}(h)
    =
    \frac{
        \pi_t^-(h)q_{t,h}(c_t)
    }{
        \sum_{u\in A_t}
        \pi_t^-(u)q_{t,u}(c_t)
    } .
    \label{eq:modelv2-oral-posterior}
\end{equation}
该 choice-conditioned distribution 与模型选择预测使用同一组 $q_{t,h}$，不引入额外的 oral-specific beta。

若口头报告编码为四维中心 $\mathbf v_t$，则将其与每个假设所选类别的原型比较：
\begin{equation}
    d_t^{\mathrm{oral}}(h)
    =
    \left\|
        \boldsymbol\mu_{h,c_t}-\mathbf v_t
    \right\|_2 .
    \label{eq:modelv2-oral-center-distance}
\end{equation}
若口头报告编码为线性区域，则在同一组预先生成的 $[0,1]^4$ Monte Carlo 点上，比较口头区域与 $R_{h,c_t}$ 的交并比。中心距离或区域重叠转换为定义在 $\mathcal H$ 上的口头假设分布后，使用 Jensen--Shannon similarity 比较模型分布与报告分布：
\begin{equation}
    \operatorname{JSSim}(p,q)
    =
    1-
    \frac{
        \tfrac12D_{\mathrm{KL}}(p\|m)
        +\tfrac12D_{\mathrm{KL}}(q\|m)
    }{\log2},
    \qquad
    m=\tfrac12(p+q).
    \label{eq:modelv2-js-similarity}
\end{equation}
由于口头报告没有进入行为参数拟合，该一致性用于检验模型潜在假设状态是否具有独立的心理解释。


\subsubsection{实现冻结与最低报告要求}

正式运行和论文报告必须满足以下要求：
\begin{enumerate}
    \item 每个条件报告完整假设数、几何构成、类别映射和固定活跃容量 $M=5$；条件 1 必须明确区分 19 个基础几何划分与 38 个带标签假设；
    \item 报告独立知觉数据文件、被试噪声模式、四维 feature 重排规则和逐试次截断抽样；保存或在保留日志的确认运行中重建 $\widetilde{\mathbf x}_t$；
    \item 报告逐试次模块顺序，并明确同试次选择只使用 $\widetilde{\mathbf x}_t$、$\pi_t^-$ 和 pre-feedback $\beta_{t,h}$；
    \item 条件 3 的类别家族映射由任务设计固定，并确认三种反馈事件互斥且完备；
    \item Profile controller 只使用过去反馈、active-set posterior entropy 和预先定义的试次进程；保存每名被试选择的 family $g_s$；
    \item 保存每个试次的 profile、active set、先验、后验、两条记忆轨道、beta 和类别概率；
    \item $w_0$ 搜索明确包含 0 与 1，并正确处理边界上的不可识别参数；
    \item $\rho$ 搜索明确包含 expectation 边界 $\rho=1$；
    \item 报告所有固定 beta 参数、数值截断常数、搜索范围、停止规则、随机种子和模拟重复数；
    \item 使用全部随机重复形成边际选择概率，并同时报告 Brier、NLL 和轨迹 CRPS；
    \item 在行为拟合完成后再评价口头报告一致性，不能用口头报告选择行为模型参数；
    \item 在正式结果生成前，对六-family参照、三-profile精简候选、条件通用 beta 更新和 $w_0$ 两个端点进行单元测试、参数恢复和全样本重新拟合；三-profile候选达到预设非劣标准后才能替代六-family参照。
    \item 将 ``有限 active set'' 与 ``精确容量 $M=5$'' 分别检验；容量敏感性分析只比较预先限定的小容量候选，不把完整假设空间作为正式心理模型。
\end{enumerate}

综上，model v2 是一个容量受限、反馈驱动的动态假设系统。物理刺激首先经独立知觉测量校准的被试特异噪声转换为知觉输入；学习者随后在任务定义的大型带标签几何空间中只维持少量活跃假设。被试特异的 feedback-gated controller 根据过去表现、后验不确定性和学习进程，在多种集合维持与更新策略之间切换；任务反馈以互斥事件似然更新证据；长短期证据通过包含两个边界的双通道记忆整合；动态 beta 调节单个假设的类别锐度；连续 concentration readout 将内部规则信念转化为外显选择。该架构保留已有拟合所支持的 controller 异质性，同时把任何进一步压缩明确置于公平重优化和非劣检验之后。
